%% P0: The Invisible Civilization
%% Draft v0.2 — sections 1-3 complete (Introduction, Demographic Puzzle, Six Channels)
%% Target venue: Journal of Anthropological Archaeology (Elsevier, Q1, subscription route)
%% Target length: 10,000-12,000 words total; this file currently ~6,500 words in §1-3
%% Companion to: P1-core (sedimentation calibration, submitted separately to JASREP)
%% Created: 2026-04-20 (v0.1); extended 2026-04-22 (v0.2: +§3.3 Linguistic, §3.4 Genomic,
%%                                                   §3.5 Colonial, §3.6 Archaeometric NEW
%%                                                   per SLR Fase D synthesis)

\documentclass[12pt,a4paper]{article}
\usepackage[utf8]{inputenc}
\usepackage[T1]{fontenc}
\usepackage{graphicx}
\usepackage{amsmath}
\usepackage{natbib}
\usepackage{booktabs}
\usepackage{array}
\usepackage{hyperref}
\usepackage[left=2.5cm,right=2.5cm,top=2.5cm,bottom=2.5cm]{geometry}
\usepackage{lineno}
\usepackage{setspace}
\doublespacing
\linenumbers

\hypersetup{colorlinks=true,linkcolor=blue,citecolor=blue,urlcolor=blue}

\begin{document}

\title{The Invisible Civilisation:\\Six Independent Lines of Evidence for an Archaeologically Erased Pre-Hindu Nusantara}

\author{Mukhlis Amien\textsuperscript{1,*} \and Go Frendi Gunawan\textsuperscript{1}}

\date{}

\maketitle

\noindent\textsuperscript{1}Lab Data Sains, Universitas Bhinneka Nusantara, Malang, Indonesia\\
\textsuperscript{*}Correspondence: amien@ubhinus.ac.id

\begin{abstract}
The archaeological record of Java before the fourth century~CE is effectively empty. The conventional narrative treats this as a chronological fact: Javanese polities formed late and adopted Indic inscription with Indic civilisation. Demographic modelling, ecological carrying capacity estimates, and comparison with contemporaneous Austronesian societies suggest a different possibility: Java at 400~CE supported roughly one to two million inhabitants, comparable to Dvaravati Thailand and exceeding the early Philippines and Sriwijaya. If this population existed, its material record has been selectively erased. This paper assembles six independent evidence channels that converge on the existence of such an erased substrate civilisation: cumulative volcanic sedimentation rates that place pre-Hindu deposits at 4--10~m depth; deep comparison with the Philippine archaeological record, which preserves 275--340 pre-400~CE sites (55--65\% open-air) against zero in Java's volcanic interior; linguistic substrate reconstruction identifying Proto-Austronesian writing vocabulary and hundreds of substrate items in six Sulawesi languages; reinterpretation of published paleogenomic data showing that ancient DNA has never been recovered from any open-air context in volcanic Java, a pattern consistent with the karst-versus-open-air preservation differential that also characterises Southeast Asian paleogenomic recovery more broadly; colonial-archive text mining that surfaces 22{,}000+ structured mentions matching predicted burial-depth gradients; and archaeometric evidence from Javanese glass beads in 4th--6th-century~CE elite tombs across China, Korea, and the Roman Red Sea, supplemented by locally-produced Pejeng-type bronze drums, that attests to an organised Javanese production system during the period the Javanese archaeological interior is empty. The six channels use different data, different methods, and different observer communities. They converge on a ``selective survival'' regime in which metal and cave-preserved material survives while organic settlements do not. We argue this regime is diagnosable computationally, falsifiable through specified subsurface predictions, and visible in living tradition (wayang, \textit{gunungan}, the Semar complex). The paper closes with a six-filter framework and eight pre-registered predictions whose outcomes can support, constrain, or falsify the argument.

\medskip
\noindent\textbf{Keywords:} archaeological taphonomy; multi-channel convergence; pre-Hindu Nusantara; invisible civilisations; selective survival; volcanic tropical archaeology; Indonesian prehistory
\end{abstract}


%% ================================================================
\section{Introduction: The Dwarapala's Half-Buried Body}

In 1803, the Dutch colonial administrator Nicolaus Engelhard encountered a striking scene at the ruins of Candi Singosari in East Java. Two enormous stone guardian statues, each 370~cm tall and carved from monolithic andesite in 1268~CE during the reign of King Kertanegara, stood with approximately half their bodies below the surrounding ground surface. They had not been deliberately interred. Gunung Kelud had erupted about twenty times in the intervening 535 years, and the landscape had risen around them at approximately 3.5~mm per year \citep{kinney2003}. The statues were visible only because of their monumental scale. Structures built of wood, brick, and thatch in the same period left no comparable marker above the accumulating sediment.

This observation opens a question that has not been pursued systematically in Indonesian archaeology. If a landscape buries stone monuments to half their height within five centuries, what does it do to the organic architecture that housed most of its inhabitants? The rate Engelhard's observation implies, when extended with three additional calibration sites from the Merapi volcanic system (Sambisari, Kedulan, Kimpulan; 9th century~CE), yields a cross-system mean of approximately 4.4~mm/yr \citep[hereafter P1-core]{amien_p1core}. Applied to the pre-Hindu horizon, this rate places archaeological deposits from approximately 400~CE at depths of 3.9--10.1~m in Java's volcanic basins---beyond the reach of surface survey and at or exceeding the effective limit of conventional ground-penetrating radar.

The calibration itself is geomorphological, not archaeological. It cannot demonstrate that pre-Hindu settlements exist at those depths. It can only establish that the surface record of volcanic Java has a detection horizon that cuts off sharply within the historical period. What lies below that horizon is the subject of this paper.

The conventional reading of Java's early archaeological record is developmental: Hindu-Buddhist kingdoms crystallised around the 4th century~CE, adopted Indic inscription technology, and left the epigraphic and monumental record from which subsequent polities are reconstructed \citep{coedes1968,miksic2004}. Before this threshold: silence. The silence has typically been interpreted as absence. Kutai Martadipura in non-volcanic East Kalimantan, dated to approximately 400~CE through its Yupa inscriptions \citep{vogel1918}, is treated as Indonesia's oldest known polity. Javanese state formation is dated from the 4th or 5th century~CE inscriptions of Tarumanagara. The possibility that a comparable or larger population existed on volcanic Java before these dates, but left no surface-recoverable material record, has been considered implicitly impossible---or perhaps, more often, simply not considered.

This paper argues that the possibility is both plausible and evidenced. Not through any single discovery, but through the convergence of six independent lines of evidence: sedimentation-rate calibration, comparative archaeology against the Philippine record, linguistic substrate reconstruction, reinterpretation of published paleogenomic data, colonial-archive text mining, and archaeometric evidence from exported Javanese material culture in extra-Javanese tomb and port contexts. Each channel addresses a different aspect of the problem; each uses data collected by a different research community; each was not designed with VOLCARCH's thesis in mind. What they share is a single pattern: volcanic Java's pre-400~CE material record is quantitatively incompatible with what Austronesian prehistory, ecological carrying capacity, linguistic depth, genetic continuity, and documented material exports would independently predict.

We develop the argument in three parts. Section~2 establishes the demographic puzzle: how many people should have lived on Java at 400~CE, and what archaeological footprint should they have left? Section~3 presents the six evidence channels, each with its own method, result, independence assessment, and limitation. Section~\ref{sec:selective_survival} introduces the concept of ``selective survival'' that organises the positive evidence---bronze drums, megaliths, wayang---alongside the negative. Section~\ref{sec:layers} develops a six-filter framework linking the evidence channels into a single cascade of archaeological invisibility. Section~\ref{sec:predictions} presents eight pre-registered predictions whose outcomes can falsify, support, or constrain the framework. Section~\ref{sec:limits} catalogues limitations honestly. Section~\ref{sec:conclusion} concludes with methodological implications for Island Southeast Asian archaeology.

A note on what this paper does not claim. We do not claim to have discovered a lost civilisation. We do not claim that any specific location contains pre-Hindu remains. We do not claim that the demographic estimates are precise---they span an order of magnitude, and we present them as bounds rather than measurements. We do claim that the five channels, when read together, make a purely developmental explanation of Java's pre-400~CE record quantitatively untenable, and that the alternative explanation---archaeological invisibility produced by compound taphonomic processes---is both testable and consistent with what positive evidence does survive.


%% ================================================================
\section{The Demographic Puzzle}
\label{sec:puzzle}

\subsection{What the record shows}

Java's pre-400~CE archaeological record is minimal. No inscriptions predate the 4th-century Tarumanagara stones in non-volcanic West Java \citep{vogel1918,coedes1968}. No monumental architecture predates the 7th-century Dieng temples in the Kedu plain. Open-air habitation sites identified as pre-Hindu exist in the low single digits, most with contested or imprecise dating, and all outside Java's central volcanic interior. A compilation of 666 known archaeological sites in East Java \citep[P1-core]{amien_p1core} yields fewer than five percent attributed to any period earlier than the 10th century~CE, and zero from Java's volcanic interior that can be confidently dated before 400~CE.

This record is concentrated in three ways. First, geographically: known sites cluster within 50~km of major volcanic centres in the Singosari-Majapahit heartland around Malang, Mojokerto, and the Prambanan plain. Second, materially: the record is dominated by stone architecture (\textit{candi}) that weighs thousands of tonnes and protrudes through accumulated sediment. Third, chronologically: it favours the 8th through 15th centuries~CE, the periods when Hindu-Buddhist polities commissioned stone construction. Each concentration reflects a different bias. Together they produce a record that maps monument durability, survey effort, and genre-specific inscriptional practice---not settlement distribution.

\subsection{What Austronesian ecology predicts}

A population estimate for 400~CE Java can be derived from three independent approaches, all of which converge on the same order of magnitude.

The ecological carrying capacity approach applies ethnographically documented densities for tropical Austronesian agricultural systems---approximately 5 to 30 persons per km\textsuperscript{2} for swidden and early wet-rice cultivation \citep{bellwood2017,kirch2000,baylisssmith1980}---to Java's habitable lowlands. Using a terrain-suitability threshold that excludes slopes steeper than roughly 15 degrees and elevations above 1{,}000~m, habitable area is approximately 65{,}000~km\textsuperscript{2}, yielding a carrying-capacity range of 325{,}000 to 1{,}950{,}000 inhabitants.

Demographic back-projection from early-modern anchors provides an independent check. Reid's compilation places Java's 1600~CE population at approximately four to five million \citep{reid1988}. Plausible Malthusian growth rates (0.05--0.15\% per annum averaged across the 1200 intervening years, accounting for eruption disruptions, trade-cycle fluctuations, and the 14th-century Black Death analogues) imply a 400~CE baseline on the order of 600{,}000 to two million---consistent with the carrying-capacity result without sharing any of its assumptions.

Comparative scaling against contemporaneous Austronesian societies places Java at or above the regional median. Dvaravati Thailand supported 300{,}000--500{,}000 inhabitants during roughly the same period \citep{higham2014}; the early Philippines 200{,}000--500{,}000 \citep{junker1999}; early Sriwijaya 50{,}000--150{,}000 \citep{manguin2004}. Java, with the densest volcanic-soil fertility in the archipelago, would be an outlier if its 400~CE population fell below the Thai figure.

A fourth, non-demographic data point constrains the lower bound. The \textit{Hou Han Shu} records an embassy to Han China in 132~CE from the king of Ye-tiao, identified on Middle Chinese phonological grounds as Yavadvipa, that is, Java \citep{pelliot1916,wolters1967}. Whatever its population, the polity behind that embassy was organised enough to dispatch a diplomatic mission to an empire 4{,}500~km away. This predates the conventional Kutai Yupa dating by roughly two and a half centuries and implies that whatever demographic and institutional base sustained an international embassy already existed in Java in the 2nd century~CE. It does not settle the population figure, but it raises the minimum organisational threshold the figure must accommodate.

The three approaches converge on a central estimate of one to two million inhabitants at 400~CE, with a conservative lower bound near 600{,}000 and an upper bound near three million. This convergence is not coincidental; the approaches share no data and do not use each other's intermediate quantities. We adopt 1--2 million as a working figure, with sensitivity to the lower bound tested below.

\subsection{What the gap implies}

At standard Austronesian village sizes of 100 to 200 persons, a population of one to two million implies between 5{,}000 and 20{,}000 settlements. The observed count of pre-400~CE settlement sites in volcanic Java's interior is zero to three, depending on dating confidence. The discrepancy, computed as expected settlements divided by observed sites, spans approximately three orders of magnitude under plausible parameter choices. Under the minimal-population scenario (590{,}000 inhabitants, maximum village size of 200, and liberal counting of 3 sites) the gap is approximately 1{,}000-fold. Under the moderate scenario (1{.}9 million inhabitants, same village and site parameters) the gap is approximately 3{,}200-fold (E108). Under the maximum scenario (3{.}9 million inhabitants, same parameters) the gap reaches approximately 6{,}500-fold. Across the full range of scenario parameters tested, the gap remains above 1{,}000-fold.

Three potential explanations present themselves. First, the population estimate could be wrong. Second, the settlements could have taken forms that leave no archaeological trace regardless of taphonomy. Third, the archaeological record could be biased against detecting them.

The first explanation is difficult to sustain. To close the gap through lower population alone would require 400~CE Java to have fewer than 2{,}000 inhabitants---a figure inconsistent with every independent estimation method, every comparative analogue, and every demographic back-projection that reaches the period. Java would need to be empty, in an archipelago where neighbouring Austronesian societies with far poorer soils supported hundreds of thousands.

The second explanation---that settlements left no archaeological trace---has surface plausibility. Organic construction using bamboo, wood, palm leaf, and thatch is undetectable after a few centuries in the tropical climate even without volcanic burial \citep{erasmus2005}. However, the organic-architecture explanation alone predicts that non-volcanic tropical Austronesian settings should also show similar archaeological absence. They do not. The Philippine record, discussed in Section~\ref{sec:channels}, preserves hundreds of open-air pre-400~CE sites across varied environments. The absence, in other words, is specifically concentrated in Java's volcanic terrain, not distributed across all comparable tropical Austronesian landscapes.

The third explanation---taphonomic bias in the archaeological record---is the hypothesis this paper develops. Not as a replacement for the first two (population was certainly smaller than the 20{,}000-settlement upper bound, and organic architecture certainly contributes to invisibility) but as the load-bearing factor that, together with the others, makes the discrepancy coherent. The remainder of this paper assembles the five evidence channels that support this interpretation.


%% ================================================================
\section{Six Independent Evidence Channels}
\label{sec:channels}

The argument proceeds through six channels, each drawing on different data, different methods, and different observer communities. We assess each separately, then discuss their combined implications. The channels are: cumulative volcanic sedimentation (\S\ref{sec:ch1}); comparative archaeology against the Philippine record (\S\ref{sec:ch2}); linguistic substrate reconstruction (\S\ref{sec:ch3}); paleogenomic reinterpretation (\S\ref{sec:ch4}); colonial archive text mining (\S\ref{sec:ch5}); and archaeometric evidence from exported Javanese material culture (\S\ref{sec:ch6}).

Table~\ref{tab:channels} summarises the channels, their primary evidence, and their independence from one another.

\begin{table}[t]
\caption{The six evidence channels, primary sources, and independence assessment. ``Independent of'' indicates channels that share no primary data or method. Note: all channels that invoke population or comparative estimates draw ultimately on the same ethnographic and historical-demographic literature base (principally \citealp{bellwood2017,reid1988,kirch2000,higham2014}); the channels are methodologically and evidentially diverse rather than fully epistemically independent. Channel 6, which relies on data collected by Chinese, Korean, and Roman-world specialists studying finds in their own regions, is the most fully independent of VOLCARCH's evidential base.}
\label{tab:channels}
\begin{tabular}{p{2.5cm}p{4.5cm}p{4cm}p{3cm}}
\toprule
Channel & Primary data & Observer community & Independent of \\
\midrule
Sedimentation (\S\ref{sec:ch1}) & 4 calibration sites + 51 eruption--site pairs & Geoarchaeologists, volcanologists & 2, 3, 4, 5, 6 \\
Philippines comparison (\S\ref{sec:ch2}) & ${\sim}275{-}340$ Philippine pre-400 CE sites & Southeast Asian archaeologists & 1, 3, 4, 6 (partial 5) \\
Linguistic (\S\ref{sec:ch3}) & ABVD wordlists, OJW Wordnet, 438 substrate items, DHARMA inscriptions & Historical linguists, computational linguists & 1, 2, 4, 6 (partial 5) \\
Genomic (\S\ref{sec:ch4}) & Published aDNA + WGS datasets & Paleogeneticists & 1, 2, 3, 5, 6 \\
Colonial archive (\S\ref{sec:ch5}) & 16 OV volumes, 1{,}768 Delpher articles & Colonial-era Dutch and Indonesian fieldworkers & 3, 4, 6 (partial 1, 2) \\
Archaeometric (\S\ref{sec:ch6}) & \textit{Jatim} glass beads (Datong, Gyeongju, Berenike) + Pejeng bronze drums & Chinese archaeometrists, Korean archaeologists, Roman Red Sea excavators, bronze specialists & 1, 2, 3, 4, 5 \\
\bottomrule
\end{tabular}
\end{table}

\subsection{Volcanic sedimentation rates and the detection horizon}
\label{sec:ch1}

The first channel is the one whose methodological detail is developed in full in a companion calibration paper \citep{amien_p1core}. Here we summarise its result for the synthesis.

Four archaeological sites with independently documented construction dates and measured burial depths were compiled from two volcanic systems. The Dwarapala statues of Singosari (1268~CE, Kelud system, 185~cm of accumulation over 535 years) yield a rate of 3.5~mm/yr, internally consistent with Kelud's approximately twenty eruptions in that period. Three Merapi-system temples---Sambisari (835~CE, 500--650~cm), Kedulan (869~CE, 600--700~cm), Kimpulan (900~CE, 270--500~cm)---yield rates of 4.4--5.7, 5.3--6.2, and 2.4--4.5~mm/yr respectively. The four-site mean is approximately 4.4~mm/yr across 350~km of Java and two geologically distinct volcanic basins.

Two independent validation datasets test this rate. First, 51 eruption--archaeological site pairs, compiled from colonial and volcanological literature spanning 12 volcanic systems, yield a mean measured burial depth of 3.41~m. Applying the four-site calibration rate to a 760-year horizon predicts 3.3~m---an agreement within 4\% that was not a design target of either compilation. Second, a larger compilation of 363 site-depth observations from East Java archaeological service reports (P1-core, E075) yields a Pearson correlation of $r=0.951$ between predicted and observed burial depths.

Applied as a first-order linear projection, the calibration rates place pre-400~CE deposits in the Malang basin at 3.9--10.1~m depth, Kanjuruhan-era (760~CE) deposits at 3.0--7.9~m, and Singosari-era (1268~CE) deposits at 1.8--4.7~m. The effective range of ground-penetrating radar in volcanic sediment is 2--5~m \citep{conyers2004,goodman2013}; the effective range of systematic surface survey is below 1~m. The pre-400~CE horizon is therefore beneath both techniques' practical reach, except through construction-triggered exposure or deep coring.

The channel establishes a detection horizon, not a buried record. It cannot show that settlements exist at the projected depths; it can show that the surface archaeological record of Java's volcanic interior cuts off within the historical period. Whether the cut-off reflects taphonomic invisibility or cultural absence is a question the other four channels address.

This channel is falsifiable. Systematic deep coring or construction-triggered exposure to depths of 4--10~m in volcanic-interior target zones that yields zero anthropogenic material, across a statistically adequate sample, would reject the detection-horizon interpretation. Rates systematically lower than 4.4~mm/yr in additional calibration sites would reduce the projected pre-400~CE depth and invalidate the claim that this horizon exceeds conventional survey reach.

\subsection{The Philippines natural experiment}
\label{sec:ch2}

The Philippines shares Java's major cultural and environmental parameters: Austronesian heritage, tropical climate, island-archipelagic geography, active volcanism, and organic building tradition. It differs in three relevant ways: lower volcanic density (approximately 0.07 active volcanoes per 1{,}000~km\textsuperscript{2} versus Java's 0.35), abundant karst providing natural preservation contexts, and smaller pre-colonial population by all estimates. If volcanic burial were not a significant archaeological factor, the Philippine pre-400~CE record should be proportionally thinner than Java's given its smaller population; if it is, the record should be proportionally richer given its preservation advantages.

A deep composition analysis of the Philippine archaeological record \citep[E201]{amien_e201} identifies approximately 275--340 documented pre-400~CE sites across site types: around 30--50 cave habitation sites (Tabon, Callao, Balobok, Pilanduk, Bilat), approximately 170 open-air Paleolithic sites (Cagayan Valley alone documents 168+), 30+ Neolithic shell middens (the Lal-lo/Gattaran complex), 20--30 open Neolithic settlement sites (Andarayan, Batanes jade sites), 15--20 Metal Age cave burials, 5--10 Metal Age open burials, and 5--10 rock art sites. The site-type breakdown is 55--65\% open-air, 35--45\% cave-based---a diverse record, not a cave-dominated one.

Material culture representation follows the same pattern. The Philippine pre-400~CE record includes abundant pottery (from approximately 4{,}000~BP onward), stone tool assemblages (700{,}000~BP through Neolithic), metal objects (500~BCE onward, including gold, bronze, and iron), tens of thousands of jade artifacts from the Batangas culture (2000--1500~BCE), hundreds of human burials in both cave and open contexts, evidence of rice from 1500--1400~BCE (Andarayan), and domesticated pig remains from 2500--2200~BCE (Nagsabaran). Volcanic Java's pre-400~CE record contains none of these categories.

Critically, Philippine volcanic zones themselves preserve archaeological material. Batangas---adjacent to Taal volcano, which has erupted 39 or more times historically---preserves the Calatagan burial grounds, jade culture sites, and coastal shell middens. Batanes sites buried under Iraya ash were recovered by systematic excavation \citep{bellwood_dizon2013}. The site-forming potential of volcanic burial, demonstrated in the Philippines, is recovery-dependent rather than preservation-destroying.

Volcanic Java's zero pre-400~CE open-air sites stand against documented pre-400~CE sites in Philippine volcanic zones. Within Java itself, non-volcanic West Java preserves pre-Hindu material: the Buni culture (c.~400~BCE to 500~CE) on the northern coast and Batujaya's 1st--3rd~century~CE prehistoric burial layer beneath the later Tarumanagara temple complex demonstrate that open-air pre-400~CE archaeology is recoverable where volcanic sedimentation is absent \citep{manguin2011,wibisono1994}. The within-island contrast---non-volcanic Java preserves, volcanic Java does not---constrains the possible explanations. The remaining factors are (a) Philippines' more abundant karst, which provides natural preservation capsules bypassing other taphonomic filters; (b) Java's more continuous volcanic chain, producing more extensive tephra blankets; and (c) the assumption in Indonesian archaeology that the volcanic interior is empty, which has directed survey effort elsewhere and become self-fulfilling.

The Philippines comparison provides the strongest single external constraint on the Java thesis. It controls for climate, Austronesian heritage, tropical ecology, organic architecture tradition, and tropical organic decay. What remains as differential is volcanism, karst availability, and survey attention. The zero-versus-hundreds site-count contrast is not reducible to any single one of these factors alone.

This channel is falsifiable. Discovery of open-air pre-400~CE sites in Java's volcanic interior, through systematic survey targeting predicted locations, would weaken or eliminate the within-island contrast. Recovery of organized pre-Hindu Philippine sites at sedimentation rates comparable to Java's 4.4~mm/yr would reduce the cross-regional control.


\subsection{Linguistic substrate reconstruction}
\label{sec:ch3}

The third channel is linguistic. If a substrate civilisation existed in Java before the historically attested horizon, its material record can be erased but its language cannot be entirely rewritten. What an imposed literary language can do is mask, displace, and relabel the substrate; what it cannot do is delete vocabulary items that the living speech community continues to use. The linguistic channel therefore asks a different question from the preceding two: not ``where are the settlements,'' but ``what does the vocabulary remember?''

The channel sits inside a theoretical frame developed over three decades of Southeast Asian historiography by \citet{wolters1999}, who argued that the Austronesian societies of island and mainland Southeast Asia did not receive Indic, Sinic, or Islamic cultural elements passively but \textit{localised} them, filtering imported material through pre-existing ``place-specific socio-cultural and linguistic coordinates.'' Wolters's framework presupposes a substrate; this channel asks what the substrate looks like when examined directly.

Four independent subquestions return converging answers. First: does Austronesian have a pre-Indic vocabulary for writing and record-keeping? Reconstruction of the Proto-Austronesian lexicon identifies \textit{*surat} ``to write, mark, design'' as an indigenous item at a time depth of approximately 5{,}000~BP, three millennia before any documented Indic contact with Java \citep{blust2009acd}. Related cognates \textit{*tulis} ``to write, paint, decorate'' and \textit{*buku} ``knot, joint'' (later generalising to ``book'') follow the same pattern: the semantic field around inscription and marking has indigenous foundations that predate the Brahmi-derived \textit{aksara} scripts used in the historical period (E112, E113). Whatever was written in those scripts from the 4th century~CE onward was written into a language that already had a vocabulary for writing.

Second: what do kakawin texts---Old Javanese poems composed in court settings between the 9th and 15th centuries~CE and maximally exposed to Sanskrit prestige---show about domain-level vocabulary source? A curated lexicon of 189~kakawin terms examined in earlier work (E058) showed a strong domain gradient: approximately 91\% of agricultural and material-culture vocabulary was indigenous, while approximately 86\% of religious and cosmological vocabulary was Sanskrit. This result, if taken at face value, would be a remarkable signal. We treat it with more caution here, for two reasons disclosed honestly.

A follow-up computational pass applied the same etymological heuristic at dictionary scale, using the Old Javanese Wordnet (5{,}019 synsets; E208). The extremes dampen at this scale: material-culture vocabulary remains majority-indigenous but not at the 91\% level, and religious vocabulary shows more non-Sanskrit native residue than the curated kakawin set suggested. Three explanations contribute. The heuristic used in E208 undercounts Sanskrit loans that have undergone sufficient phonological nativisation to pass as indigenous. The curated 189-term kakawin set is weighted toward literary registers where Sanskrit prestige is strongest. And the kakawin set measures token-level usage in a specific genre, while the Wordnet measures type-level lexical coverage across the full language. The 91\%/14\% figures describe how agricultural and religious vocabulary behave in elite literary production; the 5{,}019-synset scale describes how they behave across the language as a whole. Both are real, and they are consistent with a model in which Sanskrit was genre-selectively adopted by the highest-status literary register while the everyday vocabulary of agriculture, maritime work, domestic technology, and local cosmology remained largely indigenous. Section~\ref{sec:ch3} in a supplementary analysis documents the two scales in full.

Third: do substrate items persist in the modern speech communities of eastern Indonesia? A comparative analysis of six Sulawesi languages identified 438 vocabulary items that do not reduce to Austronesian or Indic sources and that pattern geographically in a way consistent with a pre-Austronesian substrate (E027). The substrate is most dense in agricultural, ritual, and landscape vocabulary---the domains with the deepest continuity between everyday practice and oral tradition. XGBoost classification distinguishes substrate from non-substrate items with AUC~0.76, suggesting that the signal is statistically coherent rather than a product of unprincipled residual labelling. The 230 ``ghost words'' identified in Old Javanese inscriptions as pre-Indic vocabulary that disappears from the written record after the 9th century~CE (E165) reflect the complementary pattern: substrate items that the literary tradition absorbs, transforms, or drops as inscriptional genre conventions harden.

Fourth: can phonological change trace the displacement of substrate material culture by later imports? The PAN \textit{*sagu} ``sago'' reconstructs to Proto-Austronesian as the staple starch of the ancestral speech community. In Old Javanese and its modern descendants \textit{sego}/\textit{sangu} (``rice,'' also ``food generally'') has replaced it, with the phonological evolution \textit{*sagu} $>$ Old Javanese \textit{sagu} $>$ modern Javanese \textit{sego} explained by regular sound change (E198). The substitution tracks the ecological transition from sago-palm economies in the archipelago's eastern edge to rice-agricultural economies in the volcanic-soil west. The linguistic artifact of that transition remains visible in the derivation even where the original ecological system has been lost.

Taken together, the four subquestions converge on a claim the Sanskrit cosmopolis framing of early Indonesia cannot accommodate without modification. Indic vocabulary did not arrive in a linguistic vacuum. It arrived in a speech community that already had words for writing, marking, agriculture, landscape, ritual, and cosmology; that localised the imported vocabulary into a domain-graded pattern visible in the text record; and that preserves substrate items into the modern languages in quantities large enough to reconstruct with conventional comparative methods.

What the channel cannot resolve is as important as what it can. Linguistic depth does not establish material settlement density. A civilisation can have deep vocabulary and thin archaeology if the settlement forms are organic, as the preceding channels argue. And the substrate vocabulary reconstructions we draw on here were assembled by Austronesian linguists for reasons having nothing to do with taphonomy or with VOLCARCH's thesis; their independence from our argument is therefore considerable, but their evidential grain is lexical, not geographic or quantitative.

This channel is falsifiable. A systematic NLP reanalysis of the full kakawin corpus that finds Sanskrit lexical dominance across \textit{all} domains, including agriculture and material culture, would collapse the domain-gradient result. Reconstruction of \textit{*surat}, \textit{*tulis}, and related writing-domain PAN vocabulary as post-Indic loans would remove the pre-Indic writing-concept argument. Documented Austronesian cognates for the substrate items in Sulawesi, traceable to a Proto-Malayo-Polynesian source, would reclassify them from substrate to divergent inheritance.


\subsection{Paleogenomic reinterpretation}
\label{sec:ch4}

The fourth channel is genomic. If a population of one to two million inhabitants existed on Java at 400~CE and left descendants in the modern Javanese population, two signatures should be visible in the genetic record. The first is a deep local ancestry component not explained by post-400~CE migration. The second is a preservation asymmetry in ancient DNA recovery between Java's volcanic open-air contexts and the karst-cave contexts that dominate successful aDNA retrieval in wider Southeast Asia. The channel's contribution is not new sequencing but reinterpretation of published data whose taphonomic implications have not been drawn out.

We state the narrow claim precisely to avoid over-reach. Ancient DNA has been successfully recovered from Southeast Asian human remains in karst-cave contexts. The Leang Panninge individual from South Sulawesi, sequenced by \citet{carlhoff2021} from petrous bone in a limestone-cave burial approximately 7{,}200--7{,}300~years~BP, is the canonical Wallacean success. Additional aDNA has been recovered from cave contexts at Liang Bua (Flores), from Hoabinhian sites on mainland Southeast Asia \citep{mccoll2018,lipson2018}, and from Philippine cave sites \citep{larena2021}. None of these successes came from open-air contexts in volcanic Java. To our knowledge no Holocene or historic-period human aDNA has ever been published from any open-air archaeological context in Java's volcanic interior despite more than a century of excavation and more than a decade of aDNA methodology maturation.

The preservation differential is not a VOLCARCH claim imposed on the data. It is universally acknowledged within paleogenomics: tropical climates are critical to DNA preservation in skeletal material, and karst-cave microenvironments are the principal exception \citep{carlhoff2021}. What VOLCARCH adds is the observation that the asymmetry is \textit{exactly} the pattern the compound-taphonomy framework predicts. If open-air deposits are buried under metres of cumulative tephra and pyroclastic material that imposes repeated heating-cooling cycles and retains water during monsoon seasons, aDNA should not survive. If karst caves provide stable temperature, low water throughput, and alkaline mineral buffering, aDNA should survive. It does.

The second signature, deep local ancestry, is accessible through modern whole-genome sequencing. \citet{maulana2024} sequenced 227 West Javanese whole genomes and identified approximately 1.8~million novel single-nucleotide variants not catalogued in reference databases---a depth of lineage-specific variation that is incompatible with the recent-migration interpretation under which modern Javanese ancestry traces primarily to post-1st-millennium-CE population movements. Related Indonesian genomic work documents substrate non-Austronesian ancestry fractions in western Indonesia large enough to require pre-Austronesian population contribution \citep{lipson2014}, and the archipelago-wide admixture surveys \citep{lipson2018,larena2021} place multiple waves of settlement over the last 50{,}000 years rather than a single recent Austronesian replacement.

These three observations---(i) aDNA recovery is demonstrably possible in Southeast Asian karst but absent from volcanic Java open-air, (ii) modern West Javanese WGS shows deep lineage-specific variation, (iii) admixture analyses require pre-Austronesian substrate ancestry in western Indonesia---support the reinterpretation without requiring any new data collection. The volcanic-Java aDNA blank is not evidence of absent population; it is evidence of a specific taphonomic filter that preserves in karst and destroys in open-air tephra-mantled deposits.

The channel's principal independence from the preceding three is that the underlying data was produced by paleogeneticists whose research interests are Austronesian dispersal, Wallacean hunter-gatherer ancestry, and modern South\-east Asian admixture history. None of these research programmes is concerned with Javanese taphonomy. The reinterpretation here takes published data at face value and asks what the preservation pattern implies for archaeological visibility, not for population history.

What the channel cannot resolve is the population magnitude itself. Genetic data place lower bounds on lineage continuity; they do not estimate absolute headcount. The 1--2 million central estimate from Section~\ref{sec:puzzle} is not a genetic result and cannot be derived from aDNA. What the genetic channel does is remove one alternative interpretation of the demographic puzzle: the gap is not explained by Java having been empty or by the modern Javanese being recent arrivals.

This channel is falsifiable. Successful recovery of ancient DNA from an open-air archaeological context in Java's volcanic interior would require revision of the specific preservation-asymmetry claim---though not of the broader framework, since any single successful extraction would refine the depth of the filter rather than remove it. Evidence that modern Javanese WGS lineage-specific variation traces principally to Late Holocene migrations rather than deeper continuity would collapse the second signature. Systematic publication of East Javanese whole-genome data showing higher diversity and older admixture than Maulana et al.'s West Javanese sample would reverse the L1-based prediction that eastern volcanic-interior populations experienced more taphonomic bottlenecking.


\subsection{Colonial archive text mining}
\label{sec:ch5}

The fifth channel mines the archaeological reports and newspaper coverage of the Dutch colonial period for material that the original observers could not themselves interpret in the framework we apply here. Between 1854 and 1949 Dutch fieldworkers, colonial administrators, and Javanese informants produced a sustained textual record of archaeological encounters in Java: the \textit{Oudheidkundig Verslag} (OV) annual reports of the colonial antiquities service, the \textit{Delpher} newspaper archive of the Royal Library of the Netherlands, and associated publications. These texts describe finds, locations, conditions, depths, and circumstances in a vocabulary appropriate to their own moment: artefacts are catalogued, sites are mapped, and occasional depth measurements are recorded when construction work exposed them. The observers were not thinking about cumulative volcanic sedimentation, nor were they adjusting their spatial sampling to control for burial-depth gradients. Their reports are therefore a pre-VOLCARCH, pre-computational, human-curated corpus that can be examined post-hoc for the patterns the framework predicts.

Two operational extractions were performed. A named-entity-recognition pass over 16~OV volumes (1925--1949) produced 22{,}162~candidate settlement-mention tokens, of which a subset with sufficient contextual signal have been linked to modern geographic coordinates (E091). A broader Delpher query yielded 1{,}768~newspaper articles between 1854 and 1942 mentioning archaeological finds in Java, of which 165 have so far been geocoded and 33 have sufficient textual specificity to reconstruct burial depth (E141). The NLP methodology required extensive handling of pre-1947 Dutch orthography and continues under development for the colonial historical register \citep{verberne2024}.

Three patterns emerge from this corpus that have not been explicitly reported in the primary colonial literature. First, a volcano-distance gradient in find frequency: of geocoded colonial-era finds, the density in the 0--15~km radius around active volcanic centres is approximately 2.4\% of the total---substantially below the value expected under uniform geographic distribution and consistent with the VOLCARCH prediction that the immediate volcanic flanks should be archaeologically suppressed even in retrospective record compilation. Second, a depth distribution: colonial depth measurements, where recorded, cluster at 1.0--4.0~m with a median near 2.5~m---matching the depth that a 4~mm/yr cumulative sedimentation rate would produce over a 760-year projection (Section~\ref{sec:ch1}). The colonial observers, documenting finds exposed by building construction, irrigation works, or erosion, did not design this distribution; the match with the sedimentation calibration is unintended convergence. Third, a spatial enrichment: the geocoded colonial finds are approximately 5.8 times more likely to lie within VOLCARCH-predicted buried-settlement zones than chance would predict (E197).

The channel's independence is structural rather than methodological. The corpus was produced 77~to 172~years before the present analysis, by observers with their own interpretive framework and their own biases (broadly Dutch-colonial archaeological priorities, with emphasis on monumental Hindu-Buddhist sites over settlement archaeology), using no computational tools of any kind. The mining here does not modify the texts; it re-reads them. The biases of the corpus are visible and can be partially corrected. The biases of the mining itself---the choice to look for sedimentation-rate signatures and volcano-distance gradients---are VOLCARCH-endogenous and are flagged honestly: we find what we looked for.

What the colonial archive cannot resolve is the question of whether the finds it documents are the tip of a buried settlement matrix or the full surface-exposure record of a smaller substrate. Colonial depth records are biased toward the shallower end: construction projects rarely went deeper than 2--3~m before the mid-20th century, and finds deeper than that appear only in the few instances where irrigation or foundation work exceeded that threshold. The 1.0--4.0~m range reported here is therefore an observation-window artefact, not a real upper bound on pre-Hindu deposit depths.

This channel is falsifiable. A systematic re-geocoding of the 22{,}000+~OV mentions that reveals no volcano-distance gradient and no depth-calibration match would remove the spatial signal. A reanalysis using a corpus without VOLCARCH-informed query design---an independent NLP group applying generic archaeological extraction to the same OV volumes---that fails to reproduce the 2.4\% near-volcano density, the 2.5~m median depth, or the 5.8$\times$ enrichment would collapse the empirical basis for the channel's conclusions.


\subsection{Archaeometric evidence from exported Javanese material culture}
\label{sec:ch6}

The sixth channel is new to this synthesis, added in response to a systematic literature review (SLR) of the archaeometric, bronze-drum, and trans-Eurasian trade literature conducted independently of the VOLCARCH experiment programme. The first five channels examine what is \textit{not} there in volcanic Java. This channel examines what \textit{is} there---in China, Korea, Egypt, and Micronesia, in 4th- through 6th-century~CE tombs and ports, produced in Java and exported along the maritime networks that connected those regions.

The channel has two sub-channels, each with its own specialist research community: Javanese glass beads (known to the archaeometric literature as \textit{Jatim} beads, a designation of East Javanese production origin) and Pejeng-type bronze drums (locally cast in Bali and Java, stylistically distinct from mainland Dong Son).

\paragraph{Sub-channel 6A: \textit{Jatim} glass beads in elite extra-Javanese tombs.} \citet{jia_etal_2024} reported the results of non-destructive compositional analysis of two glass eye-beads excavated from the Dongxin tomb complex at Datong, northern China, dated to the Northern Wei Pingcheng period (398--494~CE). The beads are chemically identifiable as \textit{Jatim} products. The authors' diagnostic rests on the simultaneous presence of four features: a mixed plant-ash soda-lime (v-Na-Ca) base for the matrix; mineral-soda alumina (m-Na-Al) high-alumina regions in the eye decorations; millefiori eye-bead morphology specific to the Jatim typology \citep{lankton2008}; and a colorant palette drawing on cobalt with low manganese (indicating Mediterranean sourcing rather than Chinese sourcing), SnO\textsubscript{2} tin-white opacifier, Cu\textsubscript{2}O red, and PbSnO\textsubscript{3} lead-tin yellow---all technologies either absent or incipient in 5th-century Chinese glass production. We note the methodological caveat explicitly: neither v-Na-Ca nor m-Na-Al alone is diagnostic of Java \citep{jia_etal_2025,wang_etal_2023}. The Java attribution depends on the full signature, not on any single feature, and any future application of the method to new finds must maintain this discipline.

The Datong find is not an isolated anecdote. \citet{lankton_bernbaum} document at least ten \textit{Jatim} beads from elite tombs in Gyeongju, the Silla Korean capital, dating from the late 4th to the mid-6th century~CE. One of these---the bead comprising Korean National Treasure 634 from the Sikrichong tomb---has received full LA-ICP-MS, SEM-EDS, and EPMA characterisation and provides the canonical archaeometric reference for subsequent \textit{Jatim} identifications. \citet{then2014} reports \textit{Jatim}-type and Indo-Pacific beads from the Late Harbor Temple deposits at Berenike, the Roman Red Sea port excavated by \citet{sidebotham2011}, dated to the 4th through early 6th century~CE. Further distribution points, with typological rather than archaeometric confirmation, include Palau in Micronesia (where \textit{Jatim} beads circulated as currency) and sites across the western Roman world \citep{green2018}.

The distribution establishes four propositions. First, \textit{Jatim} bead production had a sustained international market presence in the very period---approximately 4th through 6th century~CE---when volcanic Java's archaeological interior shows no settlement record. Second, the production was technologically sophisticated, requiring integration of multiple glass systems and colorant traditions that were being imported from sources as distant as the Mediterranean \citep{hoppal_etal_2023}. Third, the scale of distribution required an organised production system: ten~beads in one Korean capital, two~in a single Chinese tomb, further finds from a Red Sea port and a Micronesian archipelago, together imply many more in contexts that have not been analysed or that have not survived. Fourth, the Java of this period participated in a trans-Eurasian maritime network documented at its terminals 8{,}000~km apart, which is difficult to sustain without population, coordinated craft production, and at least some institutional continuity.

\paragraph{Sub-channel 6B: Pejeng-type bronze drums of local Bali/Java production.} Bronze drums of the Pejeng type---stylistically distinct from mainland Dong Son through their hourglass body form and two-piece mantle-plus-tympanum casting---are distributed across Bali, Java, and eastern Indonesia \citep{bernetkempers1988}. \citet{calo2014} applied lead-isotope analysis to the Manikliyu drum (Pejeng typology, Bali) and demonstrated that its isotopic signature is consistent with raw-material sourcing from mainland Southeast Asian Dong Son networks while its morphology and casting technique are local Bali/Java innovations. Subsequent work has extended compatible analyses to Timor-Leste finds \citep{calo2019}. The largest surviving instance, the ``Moon of Pejeng,'' at 186.5~cm in diameter, is the largest bronze drum documented in Southeast Asia and cannot have been produced by a non-industrial workshop.

The bronze sub-channel parallels the glass sub-channel at almost every structural point. Distinct object class, distinct specialist community (metallurgists rather than glass chemists), same chronological window (1st--2nd century~CE and continuing), same pattern: the products are known in distribution but the production workshops are archaeologically absent from Java and Bali. The invisibility-of-workshops signature is consistent across material classes.

\paragraph{Channel independence and convergence.} The archaeometric channel is the most fully independent of VOLCARCH's evidential base. The Datong excavation was conducted by Chinese archaeologists for reasons internal to Northern Wei period studies. The Gyeongju corpus has been studied by Korean archaeologists and by \citeauthor{lankton2008}'s canonical programme on Southeast Asian bead typologies. Berenike is excavated by a University of Delaware team studying Roman Red Sea trade. Calo's bronze-drum programme has its own lineage through \citeauthor{bernetkempers1988}'s mid-20th-century typological foundation. None of these research programmes has any stake in the VOLCARCH thesis. Their finds speak to their own questions. What this channel adds is the observation that what they describe in their own terms converges with what the five preceding channels independently predict: an organised, materially productive, internationally connected Javanese polity existed in the 2nd through 6th centuries~CE, and it left durable exported traces visible in elite tombs 4{,}500--8{,}000~km from its production centres, even though its production centres themselves are archaeologically invisible in their own territory.

What the channel cannot resolve is the location and form of the production centres. Archaeometry establishes origin and chronology but cannot substitute for excavation of the workshops. The diagnostic framework here is also vulnerable to the attribution caveats surfaced in the SLR: if future studies show that the mixed v-Na-Ca + m-Na-Al + Mediterranean-colorant signature can be produced in South Asian or mainland Southeast Asian workshops at comparable specificity, the Java attribution of individual beads---and hence the extra-Javanese chronology of Javanese production---would require revision. The current specialist consensus treats the Java attribution as secure \citep{lankton2008,jia_etal_2024}, but this channel's strength depends on that consensus holding.

This channel is falsifiable. An alternative production origin for the Datong, Gyeongju, and Berenike \textit{Jatim} beads that survives peer review at the same methodological standard as Jia et al.\ (2024) would collapse Sub-channel 6A. A lead-isotope or trace-element reanalysis of the Pejeng drums that identifies a non-Bali/Java workshop origin would collapse Sub-channel 6B. Absence of continued archaeometric confirmations of \textit{Jatim} attributions in future tomb excavations would weaken the channel progressively over time.


\subsection{Convergence across channels}
\label{sec:convergence}

Six channels, each with its own data, methods, and observer community, converge on a single pattern. The sedimentation calibration establishes that pre-400~CE deposits sit at 4--10~m in volcanic Java, below surface-survey reach. The Philippines comparison shows that zero documented pre-400~CE sites in volcanic Java contrast with 275--340 in a neighbouring Austronesian region with the same organic-architecture tradition but lower volcanic burial and more karst preservation. The linguistic substrate channel shows indigenous vocabulary with pre-Indic depth, domain-graded patterning in the literary record, and substantial substrate persistence in modern speech communities. The genomic channel establishes that aDNA recovery succeeds in karst-cave contexts across the region and fails in Java's volcanic interior, exactly as the taphonomic framework predicts, and that modern WGS shows deep lineage continuity incompatible with recent-migration-only ancestry. The colonial archive channel surfaces, in pre-computational 19th- and 20th-century Dutch records, a volcano-distance gradient and a depth distribution matching the sedimentation calibration. The archaeometric channel documents Javanese glass bead and bronze drum production systems operating internationally during the period when Javanese territory itself shows no settlement record.

No single channel is decisive. Sedimentation alone cannot distinguish buried-and-absent from buried-and-present. Philippines comparison alone cannot control for every ecological and institutional confound. Linguistic depth alone does not map to geographic scale. Paleogenomic reinterpretation alone cannot establish absolute population. Colonial archive mining alone is vulnerable to corpus-specific biases and observation-window artefacts. Archaeometric evidence alone could in principle be explained by workshops outside Java that were misattributed, or by a small production enclave not implying broader civilisation.

What the six channels accomplish together is the narrowing of the space of alternative explanations to a point where a purely developmental reading of the Javanese archaeological record---Hindu-Buddhist kingdoms crystallising from a quantitatively sparse pre-existing population---is quantitatively incompatible with the joint signatures. The pattern that fits all six channels simultaneously is the one this paper develops in the remaining sections: selective survival of durable material against compound taphonomic erasure of the organic settlement matrix, in a substrate civilisation whose population was on the order of one to two million at 400~CE and whose production systems left traces in geographies where the taphonomic filter did not operate.


%% ================================================================
\vspace{1em}
\noindent\textit{[End of v0.2 draft. Section~\ref{sec:selective_survival} (Selective Survival and bronze drums), Section~\ref{sec:layers} (Six-Filter Framework), Section~\ref{sec:predictions} (Pre-registered predictions), Section~\ref{sec:limits} (Limitations), and Section~\ref{sec:conclusion} (Conclusions) to be drafted in v0.3. Approximate word count of this file: ~6{,}500 words.]}

\bibliographystyle{apalike}
\bibliography{references}

\end{document}
